\documentclass[10pt,a4paper]{amsart}
\usepackage[margin=19mm]{geometry}
\usepackage{amsmath,amssymb,mathtools}
\usepackage[scr=rsfs,cal=euler]{mathalfa}
\usepackage{xfrac}
\usepackage[unicode,hidelinks,bookmarks=false]{hyperref}
\newcommand{\ie}{i.e.\ }
\newcommand{\cf}{cf.\ }
\newcommand{\resp}{respectively\ }
\newcommand{\Subsection}{\subsection}
\providecommand{\nobreakdash}{\nobreak\hbox{-}}
\setlength{\emergencystretch}{2em}
\multlinegap0pt
\usepackage{amssymb}
\usepackage[all]{xy}
\newtheorem{lemma}{Lemma}
\newtheorem{theorem}{Theorem}
\DeclareMathOperator{\lcm}{lcm}
\DeclareMathOperator{\ord}{ord}
\title{On a family of non-Volterra quadratic operators acting on a simplex}
\author{Uygun Jamilov, Manuel Ladra}
\date{Source: arXiv:2601.17494v1 (CC BY 4.0)}
\begin{document}
\maketitle

\noindent\emph{Source and licence.} On a family of non-Volterra quadratic operators acting on a simplex. Version arXiv:2601.17494v1 (CC BY 4.0), distributed by arXiv under the Creative Commons Attribution 4.0 International licence, http://creativecommons.org/licenses/by/4.0/, as recorded in the arXiv licence metadata for this exact version and shown on the abstract page https://arxiv.org/abs/2601.17494v1. Full licence notice: \texttt{licenses/arxiv-2601.17494-CC-BY-4.0.txt}.\par\smallskip

\noindent\textit{Editorial scope.} Counted unit: the article's theorem theorem (source0.tex lines 801--806). The article's complete original proof of it is delivered with it (source0.tex lines 808--881, one \texttt{proof} environment of 74 lines). No mathematics is written, repaired or re-derived: the statement and the proof are the article's own lines, in the article's own notation, and no retained line is substituted or re-flowed.  Retained uncounted as the source-local context the proof needs: lines 713--721; lines 779--781.  Nothing was omitted from any retained span, so the package carries no omission record.  No retained line contains a citation, so the wrapper carries no bibliography and no bibliography entry.  No retained line contains a figure environment, an image-inclusion command or drawing code, so no image is delivered and the package declares no asset dependency.  Editorial formatting only, in addition: an AMS \texttt{amsart} preamble that restates the article's own declarations and macros used by the retained lines, namely the environments \texttt{lemma}, \texttt{theorem} on separate counters, together with the standard packages those lines need.  The retained environments print in this document as Lemma 1, Theorem 1 in order of appearance, whereas the article prints the same items with the numbering of its own full document.  The complete native source of the article as distributed in the arXiv e-print for this exact version is bundled unchanged as \texttt{sources/193317/source0.tex}.
\begin{equation}\label{opcc}
V_\alpha: \left\{\begin{array}{lll}
               x'_{k}= \alpha x_k^2 +
               \frac{2\alpha}{m-2}\displaystyle\sum_{\substack{i,j\in E\setminus\{k\}\\ i<j}} x_ix_j + 2(1-\alpha)x_m x_{\pi(k)}, \ \ 1\leq k\leq m-1\\[5mm]
               x'_{m}= x_m^2+
               \frac{2\alpha}{m-2}\displaystyle\sum_{\substack{i,j=1\\ i<j}}^{m-1} x_ix_j +
               (1-\alpha) \Big(\displaystyle\sum_{\substack{i=1}}^{m-1} x_i\Big)^2.
\end{array}\right.
\end{equation}

\begin{lemma}\label{falfa} If $0<x<1$, then
\[\lim\limits_{n\rightarrow\infty} f^n_\alpha(x)=x_m^*.\]
\end{lemma}

\begin{theorem} For the QSO $V_\alpha$ the following statements are true:
\begin{itemize}
\item[i)] The function $\phi(\mathbf{x}) =x_m$ is a Lyapunov function for the operator $V_\alpha$;
\item[ii)] The set of limit points is $\omega_{V_\alpha}\big(\mathbf{x}^{(0)}\big)=\{\mathbf{x}^*\}$ for any $\mathbf{x}^{(0)}\in S^{m-1}$.
\end{itemize}
\end{theorem}

\begin{proof}  i) Evidently the function $\phi(V_\alpha(\mathbf{x}))=x'_m$ is continuous on $S^{m-1}$ and
\[\min\limits_{\mathbf{x}\in S^{m-1}} \phi(V_\alpha(\mathbf{x})) =x_m^* \ \ \text{iff} \ \  \mathbf{x}=\mathbf{x}^*, \ \ \text{where}\]
\[x_1^*=\cdots=x_{m-1}^* =\frac{1}{ (2-\alpha)(m-1)+\alpha}, \ \ x_m^*=\frac{(1-\alpha)(m-1)+\alpha}{ (2-\alpha)(m-1)+\alpha}.\]

Using the well-known Maclaurin's inequality, we have
\[\frac{2\alpha}{m-2} \displaystyle\sum_{\substack{i,j=1\\ i<j}}^{m-1} x_ix_j =
(m-1)\alpha \frac{\displaystyle\sum_{\substack{i,j=1\\ i<j}}^{m-1} x_ix_j}{
\binom{2}{m-1}} \leq \frac{\alpha}{ m-1} (x_1+\dots+x_{m-1})^2.\]

Using the last and \eqref{opcc} one has that
\begin{align*}
\phi(V_\alpha(\mathbf{x}))-\phi(\mathbf{x})&=x'_m-x_m= x_m^2+
               \frac{2\alpha}{m-2}\displaystyle\sum_{\substack{i,j=1\\ i<j}}^{m-1} x_ix_j +
               (1-\alpha) \Big(\displaystyle\sum_{\substack{i=1}}^{m-1} x_i\Big)^2-x_m\\
&= x_m^2+(1-\alpha) (1-x_m)^2 +  \frac{2\alpha}{m-2}\displaystyle\sum_{\substack{i,j=1\\ i<j}}^{m-1} x_ix_j-x_m\\
&\leq x_m^2+(1-\alpha) (1-x_m)^2 +  \frac{\alpha}{m-1}(1-x_m)^2-x_m\\
&=x_m^2+\Big(1-\frac{(m-2)\alpha}{m-1}\Big) (1-x_m)^2-x_m\\
&= f_\alpha(x_m)-x_m,
\end{align*}

Some analysis gives us
\[\max  \Big(f_\alpha(x_m)-x_m\Big) =0, \ \ x_m\in \Big[x^*_m,1\Big]. \]


Thus $\phi(V_\alpha(\mathbf{x}))\leq \phi(\mathbf{x}) $ for all $\mathbf{x}\in S^{m-1}$, that is
the function $\phi(\mathbf{x})$ is a Lyapunov function for the operator $V_\alpha$.

Consequently, it  holds
\[\phi(\mathbf{x}^{(n+1)}) \leq \phi(\mathbf{x}^{(n)}), \quad n=0,1,2,\dots   \]
and  due to Lemma \ref{falfa}  we have
\begin{equation}\label{mcorr}
x_m^* \leq \lim\limits_{n\rightarrow\infty} \phi(\mathbf{x}^{(n+1)})
\leq \lim\limits_{n\rightarrow\infty} f^n_\alpha(x_m)=x_m^*.
\end{equation}


ii) Let $\mathbf{x}^{(0)}\in S^{m-1}\setminus \{\mathbf{e}_m\}$. Due to claim of  part i) we have $ \lim\limits_{n\rightarrow\infty} x_m^{(n)}=x_m^*<1/2$. So for enough large $n$ it  holds
\[\mathbf{x}^{(n)}\in S^{m-1}_{1/2}=\{\mathbf{x}^{(n)}\in S^{m-1}: x_m^{(n)}<\frac{1}{2}\}. \]



Using induction for $u\neq v,  \ u,v=1,\dots,m-1$, and $s=\lcm \big(\ord (\tau_1),\dots, \ord (\tau_q)\big)$, one has
\begin{align*}
x_u^{(s)}-x_v^{(s)}&= \big(x_u-x_v\big)\Big(\Big(\frac{\alpha}{ m-2}\Big)^s\prod\limits_{j=0}^{s-1} \big(m(x_u^{(j)}+x_v^{(j)}) - 2\big)
    + 2^s(1-\alpha)\prod\limits_{j=0}^{s-1}x_m^{(j)} \Big) \\
\end{align*}


For $\mathbf{x}\in S^{m-1}_{1/2}$ we have
\begin{align*}
\big|x_u^{(s)}-x_v^{(s)}\big|&=\big|x_u-x_v\big|\bigg| \Big(\frac{\alpha}{m-2}\Big)^s\prod\limits_{j=0}^{s-1} \big(m(x_u^{(j)}+x_v^{(j)}) - 2\big)
    + 2^s(1-\alpha)\prod\limits_{j=0}^{s-1}x_m^{(j)}\bigg|\\
    &\leq \big|x_u-x_v\big| \bigg| \Big(\frac{\alpha}{m-2}\Big)^s\prod\limits_{j=0}^{s-1} \big|m(x_u^{(j)}+x_v^{(j)}) - 2\big|+
     (1-\alpha)\bigg|\\
    &\leq  \big|x_u-x_v\big| \big(1-\alpha+\alpha^s\big).
\end{align*}
Therefore for $u\neq v, \  u,v=1,\dots,m-1$, it follows
\begin{equation}\label{1coor}
\lim\limits_{n\rightarrow\infty} \big|x_u^{(s\cdot n)}-x_v^{(s\cdot n)}\big|\leq \big|x_u-x_v\big| \lim\limits_{n\rightarrow\infty} \big(1-\alpha+\alpha^s\big)^n=0.
\end{equation}

Consequently, from \eqref{mcorr} and \eqref{1coor}, it follows that for the operator $V_\alpha^s$ the trajectory  converges to
the fixed point $\mathbf{x}^*$.

Let $\mathbf{x}^{(0)}\in S^{m-1}\setminus \{\mathbf{e}_m\}$ be an initial point and the sequence  $\big\{\mathbf{x}^{(n)}\big\}$ its trajectory.
As shown in above if $n=s\cdot k$ the subsequence  $\{\mathbf{x}^{(k\cdot s)}\}$ converges to $\mathbf{x}^*$  when $k\rightarrow\infty$.

Since the point $\mathbf{x}^*$ is a fixed point and the operator $V_\alpha$ is continuous
for the cases $n=s\cdot k+l, \  l=1,\dots,s-1$, we also have
\[\lim\limits_{n\rightarrow\infty} V^l\big(\mathbf{x}^{(k\cdot s)}\big) = \mathbf{x}^*.\]

Thus, the trajectory starting at any point $\mathbf{x}^{(0)}\in S^{m-1}$ converges to
the fixed point $\mathbf{x}^*$.
\end{proof}

\end{document}
